\documentclass[11pt,a4paper]{ltjsarticle}
\usepackage[margin=23mm]{geometry}
\usepackage{amsmath,amssymb,amsthm,mathtools,booktabs}
\usepackage[haranoaji]{luatexja-preset}
\usepackage[unicode,hidelinks]{hyperref}
\newtheorem{theorem}{定理}[section]
\newtheorem{proposition}[theorem]{命題}
\newtheorem{lemma}[theorem]{補題}
\newtheorem{corollary}[theorem]{系}
\theoremstyle{definition}\newtheorem{remark}[theorem]{注意}
\newtheorem{definition}[theorem]{定義}
\renewcommand{\proofname}{証明}
\newcommand{\PP}{\mathbb P}\newcommand{\kk}{\mathbb C}
\newcommand{\Sym}{\operatorname{Sym}}\newcommand{\qgr}{\operatorname{qgr}}
\newcommand{\Ext}{\operatorname{Ext}}\newcommand{\Extop}{\operatorname{Ext}}
\newcommand{\Hom}{\operatorname{Hom}}\newcommand{\HH}{\operatorname{HH}}
\newcommand{\Spec}{\operatorname{Spec}}\newcommand{\PGL}{\operatorname{PGL}}
\newcommand{\OO}{\mathcal O}\newcommand{\rank}{\operatorname{rank}}
\title{非可換三次曲面と非可換射影三空間のモジュライ\\\large 四周期表示・明示的成分・反標準復元とWeyl商}
\author{植田一石}
\date{2026年11月17日 / surface-v4}
\begin{document}\maketitle
\begin{abstract}
非可換三次曲面の反標準偏極による内在的記述を、
四周期の二次関係を持つ非可換射影三空間および四次superpotentialのモジュライと比較する。
任意に指定した四周期構造と非捩れpotential構造が一般には一致しないことを例示する。
中心一次拡大から得られる$R(1,3)$族では、四頂点potentialの全テンソル方向の
接空間を有理数上で計算し、67次元既約成分の滑らかな一般開部分と
7次元Deligne--Mumford商stackのエタール表示を得る。
滑らかな三次曲面のPoisson構造のambientへの一意延長と、
固定ambientの中心三次超曲面のペンシルを記述する。
反標準orbit環による超曲面認識部分の比較射を定義し、
可換点の形式近傍では通常の幾何学的markingの忘却を$W(E_6)$商として記述する。
AOUの一般位置における反標準模型の比較を仮定すると、
marked曲面から$R(1,3)$成分への有理ambient射は支配的であり、
一般ファイバーは一次元である。
中心三次ペンシルの特異値をJacobian環上のモニック8次多項式で記述し、
その8点を非可換ambientの$t\ne0$のpoint moduleと同定する。
一般には各臨界値にただ一つのnodalな特異曲面が対応する。
全weakly markedモジュライの代数性、全対象の超曲面認識、大域Weyl商は
ここでは証明しない。
\end{abstract}
{\small\tableofcontents}
\clearpage


\section{先行研究と対象の区別}
非可換三次曲面には、反標準偏極による定義、非可換平面の6点blow-up、
4次元AS正則代数の正規三次商、例外列からの箙関係という複数の表示がある。
これらの全族・全境界での同値は、本稿の前提にはしない。
Van den Bergh\cite[Chapter 12]{VdB}は、楕円的な非可換blow-upに付随する反標準型環を
4次元AS正則環の正規三次商として構成した。ただし反標準モデルとblow-upそのものの
区別に注意する（同章Remark 12.0.5）。

Abdelgadir--Okawa--Ueda\cite{AOU}は9頂点18矢の箙と9個の二次関係を用い、
marked非可換三次曲面のトーリックGITモジュライを構成した。その稠密トーラスは
$(\mathbb G_m)^8$である（\S5.4）。同論文Theorem 1.1は楕円decupleによる
双有理表示、Theorem 9.4は所定のblow-upとの導来同値を与える。
本稿ではこれを「反標準環の族とそのambientへの射の全域での同定」と読み替えない。
2025年のOkawa--Ueda\cite{OU2025}によるmarked AS曲面の一般化も一次資料として比較する。

Pym\cite[\S3.4]{Pym}の$R(1,3)$成分は、三変数Sklyanin代数の中心一次拡大に対応する。
中心拡大の代数構造はLe Bruyn--Smith--Van den Bergh\cite{LBS}、
非斉次potentialの中心元と量子del Pezzo構成はEtingof--Ginzburg\cite{EG}に属する。
本稿の寄与は、四窓全体の接空間を計算するとともに、
周期の選択およびmarking忘却の効果を明示して、特定の一般開部分における
対応射・ファイバー・形式変形を比較することである。
この結果の既知の記述との優先権の照合はなお必要である。

\subsection*{主結果と適用範囲}
命題\ref{prop:period-counterexample}は四周期の関係式stackと非捩れpotential表示の相違を示す。
定理\ref{thm:r13-open}は標準AS条件を満たす$R(1,3)$族の一般開部分を構成する。
定理\ref{thm:extension}、定理\ref{thm:pencil-critical}と命題\ref{prop:center}はPoisson延長と中心三次ペンシルの特異値を記述する。
定理\ref{thm:formal-weyl}は可換点の形式変形に限定される。
定理\ref{thm:aou-dominance}は一般位置の比較仮定の下での結論であり、
全非可換対象や全境界での支配性を主張しない。
以下では選択周期、構造対象、頂点スカラーのrigidificationを区別する。

\section{関係式・potential・選んだ周期構造}\label{sec:period-stack}
$V_i=\kk^4$ ($i\in\mathbb Z/4$) を頂点$i$から$i+1$への矢空間とし、
$G=\prod_i\mathrm{GL}(V_i)$、$\overline G=\prod_i\mathrm{PGL}(V_i)$とおく。
標準AS条件とは左右Koszul解消の型$(1,4,6,4,1)$、Gorenstein parameter 4、
行Hilbert級数$(1-t)^{-4}$を意味する。族では各成分の局所自由性を要求する。

\begin{proposition}[任意の周期を保持する関係式模型]\label{prop:relation-period}
$\prod_{i=0}^3\operatorname{Gr}(6,V_i\otimes V_{i+1})$内の標準AS部分関手を
$\mathcal R^{\rm AS}$とする。商stack$[\mathcal R^{\rm AS}/G]$は、
4周期同型$\tau:C\xrightarrow{\sim}C(4)$を\emph{指定した}標準AS正則
$\mathbb Z$-代数のstackと同値である。
\end{proposition}
\begin{proof}
$C_{i,i+1}$ ($0\le i<4$)の基底を選び、$\tau$によって全添字へ延長する。
二次関係の核が$R_{i+4}=R_i$を満たすので四つのGrassmannianの点となる。
逆に四つの関係を整数列上へ反復し、自由経路代数をその二次イデアルで割れば$C$を得る。
$\tau$を保つ同型は四つの基底変更で尽くされる。相対的には矢空間をfpqc局所で自明化して
同じ構成を行い、ベクトル束・積・関係の降下を用いる。AS部分への制限も両方向で一致する。
これはAS locus全体が有限型代数的stackであるという主張ではない。
\end{proof}

\paragraph{符号と微分の規約.}
$\mathcal W=V_0\otimes V_1\otimes V_2\otimes V_3$とし、最初の矢を最後へ送る回転を$c$とする。
$w\in\mathcal W$から、異なる始点成分を持つ
$W(w)=w-cw+c^2w-c^3w$を作る。これは$cW=-W$を満たす四次superpotentialである。
代数は長さ2の経路微分で$D(W,2)=\kk Q/(\partial_pW:|p|=2)$と定義する。
BSW\cite[Theorems 6.2, 6.8]{BSW}はKoszul (twisted) CY代数と、対応する
解消複体が完全な (twisted) superpotentialの対応を与える。
通常の符号なしcyclic quotientへ$W$を送ると零になるので、その意味のpotentialと混同しない。

\begin{proposition}[任意の4周期には非捩れpotentialは対応しない]\label{prop:period-counterexample}
命題\ref{prop:relation-period}の対象全体を、非捩れ四次superpotentialのモジュライと
無条件に同一視することはできない。
\end{proposition}
\begin{proof}
$C=\check{\kk[x_0,x_1,x_2,x_3]}$とし、$\sigma=\operatorname{diag}(2,1,1,1)$から
$\tau=\check\sigma\circ(\text{添字の4移動})$を選ぶ。
最初の四つの矢空間を通常の$V$と同定し、その後の基底を$\tau$で運ぶと
\[
 R_0=R_1=R_2=\wedge^2V,\qquad R_3=(1\otimes\sigma^{-1})\wedge^2V.
\]
この最後の空間は$\wedge^2V$でない。例えば$x_0\otimes x_1-\frac12x_1\otimes x_0$
が属する。
もし同じ四窓が非捩れpotentialを持てば、始点0の最高syzygyテンソルは
\[
 (\wedge^2V\otimes V^{\otimes2})\cap
 (V\otimes\wedge^2V\otimes V)\cap
 (V^{\otimes2}\otimes\wedge^2V)=\wedge^4V
\]
に属する。従って非零スカラーを除いて交代テンソルであり、その巡回微分が与える
$R_3$も$\wedge^2V$となって矛盾する。標準AS性自体は元の多項式$\mathbb Z$-代数から保証される。
\end{proof}

正しい比較は、非捩れpotentialに対してはFrobeniusに適合する周期を用いること、
任意の周期を扱う場合はNakayamaの捩れとtraceの変更も含むtwisted potentialを用いることである。
幾何学的点で適合周期の存在を示すだけでは、任意の\emph{選択}を分類したことにならない。
また$[U/\mathrm{GL}_4^4]$と$[U/\mathrm{PGL}_4^4]$は異なる。
射影potentialでは$(\mathbb G_m)^4$が全点を固定し、前者にはその慣性が残る。
本稿で7次元というときは後者、または対応する一般有理商を指す。

\section{R(1,3)の四頂点表示と有限型の一般開部分}\label{sec:r13-explicit}
四つの矢空間を一時的に$V=\langle t,x,y,z\rangle$へ同定する。
完全対称化を$\operatorname{sym}$、四変数の完全交代テンソルを$\operatorname{Alt}(txyz)$とする。
\[
 F=f_3(x,y,z)+tq_2(x,y,z)+t^2\ell_1(x,y,z)
 \quad\bmod\kk t^3
\]
は固定した中心直線に対して19係数を持つ。次のテンソルを定める。
\begin{equation}\label{eq:r13-potential}
 w_F=\operatorname{Alt}(txyz)
       -3(1-c+c^2-c^3)\bigl(t\operatorname{sym}(F)\bigr).
\end{equation}
式中の回転は、いったん全slotを$V$に同定して行う。得られたテンソルの各slotを
$V_0,V_1,V_2,V_3$へ戻して$\mathcal W$の元とみなす。
その四頂点superpotentialは上の$W(w_F)$である。

\begin{proposition}\label{prop:potential-presentation}
式\eqref{eq:r13-potential}の長さ2の微分が生成する関係空間は各始点で
\begin{equation}\label{eq:r13-relations}
 [t,x]=[t,y]=[t,z]=0,\qquad
 [y,z]=\operatorname{sym}F_x,\quad
 [z,x]=\operatorname{sym}F_y,\quad
 [x,y]=\operatorname{sym}F_z
\end{equation}
の6次元空間に等しい。$t^3$を$F$へ加えても$w_F$は変わらない。
$f_3$の平面三次が滑らかな範囲では、この族はnoetherian、Koszul、標準Hilbert型の
4次元CY代数を与える。
\end{proposition}
\begin{proof}
全二階微分を$V^{\otimes2}$の語で展開する。
$(\partial_{xy}-\partial_{yx})w_F=2[t,z]$とその巡回式から三つの中心関係を得る。
$\partial_{tx}w_F$などは、これらの中心関係で割った後に
$[y,z]-\operatorname{sym}F_x$などとなる。
逆包含は、各語$x_ix_j$を可換二次式$x_ix_j+\Pi_{ij}/2$へ写すことで直接確認できる。
ここで$\Pi_{xy}=F_z$など、$\Pi_{ti}=0$である。
この写像の核は式\eqref{eq:r13-relations}そのものである。
各式は$F$の19係数についての多項式恒等式であり、付属コードは独立な19文字を残して展開する。
$F=t^3$に対する補正は$(1-c+c^2-c^3)t^4=0$である。

AS条件について、$t$を係数として三変数potentialの低次摂動とみなす。
最高次代数はMori--Smith\cite[Theorem 4.2(1)]{MS}の$A_{f_3,1}$であり、
滑らかな$f_3$では3次元CYである。低次cyclic potentialのPBW定理
\cite[Proposition 7.3.1(ii)]{EG}により、active次数の随伴は$A_{f_3,1}\otimes\kk[t]$。
各$t$の特殊化のPBW基底を持ち上げると$\kk[t]$上自由であり、$t$は中心正則一次元となる。
中心一次拡大のKoszul AS上昇\cite{LBS}から標準4次元AS正則性を得る。
式\eqref{eq:r13-potential}と完全Koszul解消を使い、BSWによりCY性を得る。
パラメータ基底上では各固定次数の表示行列のrankが一定で、被約な基底上の局所自由性と
解消の各内部次数の完全性が従う。必要なら有限個の基底chartへ制限する。
\end{proof}

\subsection{七係数chartと全四頂点の接空間}
\begin{equation}\label{eq:seven-parameter}
 F_{\boldsymbol a}=a(x^3+y^3+z^3)+bxyz
       +t(x^2+c_2y^2+c_3z^2)
       +t^2(d_1x+d_2y+d_3z)
\end{equation}
を用いる。パラメータは$(a,b,c_2,c_3,d_1,d_2,d_3)$の7個である。
可逆なHesse変換とshearによる一般形としても理解できるが、以下では厳密な横断行列で
支配性を示すので、一意な大域正規形を仮定しない。

$Z_6\subset\PP(\mathcal W)=\PP^{255}$を二つのcyclic flattening
$M_{01|23},M_{12|30}$ ($16\times16$) がともにrank6となる局所閉schemeとする。
rank条件は全7次minorの消滅と適切な6次minorの非消滅で課す。
対応する二次関係は、これらおよび転置flatteningの像である。

\begin{theorem}[四頂点成分の一般開部分]\label{thm:r13-open}
$\overline G=(\mathrm{PGL}_4)^4$とする。
$\overline G\cdot[w_F]$の像の既約閉包$\mathcal Z_{13}$は$Z_6$の67次元既約成分であり、
一般点で正則である。ある非空開集合$N\subset\mathbb A^7$と
$\overline G$-不変な非空開部分$O_{13}\subset\mathcal Z_{13}$が存在して、
\[
 \overline G\times N\longrightarrow O_{13},\qquad(g,\boldsymbol a)\longmapsto g[w_{F_{\boldsymbol a}}]
\]
はエタール全射となる。$O_{13}$の全代数は標準Koszul AS正則である。
従って$[O_{13}/\overline G]$は7次元の滑らかな有限型DM stackで、$N$がエタールatlasである。
\end{theorem}
\begin{proof}
式\eqref{eq:seven-parameter}の点
\[
 \boldsymbol a_*=(3,5,2,3,1,2,4)
\]
を取る。最高三次のHesse判別条件は$9^3+5^3=854\ne0$であり、前命題が適用できる。
この点のflatteningを$M$とするとrankは6。
rank条件の一次式は、左核$L$と右核$K$の各組に対する
$L^{\mathsf T}\dot M K=0$である。両flatteningについて合計200本、未知数256個の行列を$T$とする。
四つの$\mathfrak{gl}_4$による軌道微分を$O$、7係数変分を加えた行列を$S_7$とする。
付属の有理数上の行列証明書により
\begin{gather*}
 \operatorname{rank}T=188,\qquad
 \operatorname{rank}O=61,\qquad
 \operatorname{rank}S_7=68,\qquad TS_7=0.
\end{gather*}
従ってprojective接空間は$256-188-1=67$次元であり、族の像の次元も少なくとも$68-1=67$。
源$\overline G\times\mathbb A^7$の次元が$60+7=67$なので、局所次元と接空間次元が一致する。
従って$Z_6$はこの点で正則であり、像の閉包はその既約成分である。
射影軌道は$61-1=60$次元なので、stabilizerのLie代数は零。
両方が滑らかな67次元空間となる近傍へ制限すると、全射な接写像は同型となり、射はエタールとなる。
そのエタール locus の像は開である。源の$N$を縮小し、その$\overline G$飽和を$O_{13}$とする。
源の全点のAS性は前命題で保証されている。相対的な次数成分と解消はエタール忠実平坦降下で
$O_{13}$へ降りる。stabilizerのLie代数零条件も開に制限できるため商はDMである。
\end{proof}

これは単にPymの一頂点7次元族の係数を数えた議論ではない。
\emph{全256テンソル方向}で接空間を計算し、異なる4周期表示の方向も含めて成分性を認定している。
ただし閉包の全境界におけるAS条件、正規性、正規化、有理不変量の生成元は未決定である。
具体的なエタール群oidは
\[
 \mathcal G_7=\{(\boldsymbol a,\boldsymbol a',g):
        [w_{F_{\boldsymbol a'}}]=g[w_{F_{\boldsymbol a}}]\}
       \rightrightarrows N
\]
であり、行列のliftでは$(G_0\otimes\cdots\otimes G_3)w_{F_{\boldsymbol a}}
=u w_{F_{\boldsymbol a'}}$、$u\ne0$、$\det G_i\ne0$という有限式になる。
一般次数や有限stabilizerの位数をまだ算出していないので、$N$を粗モジュライそのものとは呼ばない。
GL商のstack次元は$67-64=3$で、4次元スカラー慣性をrigidifyすると7となる。


\subsection{中心一次拡大の次数付き同型とアフィン商}
\label{sec:r13-affine-quotient}
以下では四頂点の任意周期の同型ではなく、\emph{一頂点の次数付き代数}
$B_f$の同型を考える。これは大域Zhang同値の主張とは異なる。
$U=\kk^3$、$\mathcal F_{19}=\{f\in\kk[U]_{\le3}:f(0)=0\}\simeq\mathbb A^{19}$
とし、$f_3$が滑らかな平面三次を定める非空開集合を
$\mathcal F_{19}^{\rm sm}$とする。
$f=f_3+f_2+f_1$を$t$で斉次化した$F_f$に対して
式\eqref{eq:r13-relations}で定まる次数付き代数を$B_f$と書く。

\begin{theorem}[R(1,3)中心拡大の次数付き同型]\label{thm:r13-affine}
$f,f'\in\mathcal F_{19}^{\rm sm}$とする。
$B_f$の次数一次の中心は$\kk t$であり、
$B_f$と$B_{f'}$が次数付き代数として同型である必要十分条件は、
ある$(A,b)\in\operatorname{GL}_3\ltimes U=\operatorname{Aff}_3$
によって
\begin{equation}\label{eq:affine-action}
 f'(u)=\rho_{A,b}(f)(u):=
 \frac{f(Au+b)-f(b)}{\det A}
\end{equation}
となることである（同型の向きを反転すれば群作用の逆を用いる）。
任意の次数付き同型は、全次数一次生成元の共通スカラー倍を除けば
$t\mapsto t$、$u\mapsto Au+b t$の形に一意に表される。
従って次数の共通$\mathbb G_m$-スカラーをrigidifyした
\emph{次数付き表示族}の同型群oidは
\[
 [\mathcal F_{19}^{\rm sm}/\operatorname{Aff}_3]
\]
で表される。
定理\ref{thm:r13-open}の有理数証明点を含むある非空不変開部分では
stabilizerが有限であり、この商は滑らかな7次元Deligne--Mumford stackである。
\end{theorem}
\begin{proof}
$B_f$の二次関係は、$\wedge^2\kk^4$から$\operatorname{Sym}^2\kk^4$
へのgraphをなすので、二次部分は完全対称化により
$\operatorname{Sym}^2\kk^4$と線型同型である。
一次元$\ell=a_1x+a_2y+a_3z+ct$が中心なら、
$\ell$と$x,y,z$の交換子はこの対称化表示で
$a\times\nabla f_3$を最高active次数に持つ。
$a\ne0$ならこの三つが零であることは、
適当な線型座標で$f_3$が一変数の三乗になることを強制し、
平面三次の滑らかさに矛盾する。
従って$Z(B_f)_1=\kk t$である。
同じ計算は族でも成り立つ。一次中心を定める有限自由加群間の
交換子写像は、$\kk^3=\langle x,y,z\rangle$上で各幾何学的ファイバーに
階数3を持つ。よって基底上で局所的に三次小行列式を可逆にでき、
一次中心は$\OO_Tt$という直線束であり、任意の底変更と可換する。
従って以下の同型の記述はfpqc局所的な族についても適用できる。

次数付き同型$\psi:B_f\to B_{f'}$は$\psi(t)=c t'$を満たす。
全次数一次元を同じ$c^{-1}$倍する次数付き自己同型を
ターゲットに合成して$\psi(t)=t'$に規格化できる。
このとき$\psi(u)=Au'+b t'$で、$A$は可逆である。
$B_{f'}$の交換子を使えば、列ベクトル表示で
\[
 \begin{pmatrix}[\psi(y),\psi(z)]\\
                 [\psi(z),\psi(x)]\\
                 [\psi(x),\psi(y)]\end{pmatrix}
 =(\det A)A^{-T}\,\operatorname{sym}(\nabla F_{f'}).
\]
一方、$B_f$の関係を$\psi$で送ると、右辺は
$\operatorname{sym}(\nabla F_f(Au'+bt',t'))$である。
完全対称化は線型変換に両立するので、
$\nabla_{u'}F_f(Au'+bt',t')=(\det A)\nabla_{u'}F_{f'}(u',t')$
が必要十分である。
従って両斉次三次式は$t'^3$の倍数を除いて等しく、
$t'=1$で定数項を零へ規格化すると式\eqref{eq:affine-action}となる。
逆に同式のもとでは六つの二次関係が全て保たれるので、
$u\mapsto Au'+bt'$は次数付き同型を与える。
これにより同型の全てと共通スカラーによるrigidificationを得る。

$\dim\mathcal F_{19}^{\rm sm}=19$、
$\dim\operatorname{Aff}_3=12$である。
一般stabilizerが有限であることは次の簡単な点で直接検査できる。
\[
 f_*=x^3+y^3+z^3+x^2+2y^2+3z^2+x+2y+4z.
\]
$\operatorname{Aff}_3$のLie代数元$(M,\beta)$による
$f$の一次変化は、定数項を除いて
\[
 Df(u)(Mu+\beta)-(\operatorname{tr}M)f(u)
\]
である。
$f_*$の三次部分でこれが零なら、
$x_i^2x_j$の係数から$M$の非対角成分は零、
$x_i^3$の係数から$M=s\,\mathrm{id}$となる。
二次部分の係数から
$\beta=(s/3,2s/3,s)$であり、一次部分の$x$係数は
$-2s+2s/3=-4s/3$となる。
従って$s=0$、$M=\beta=0$である。
この点のstabilizerは標数零の有限型代数群でLie代数が零なので有限。
有限stabilizer条件の非空不変開へ縮小し、
滑らかな商の次元$19-12=7$を得る。
\end{proof}
\begin{remark}
Pym\cite[\S3.4]{Pym}および
Le Bruyn--Smith--Van den Bergh\cite[\S3]{LBS}の
中心拡大の一般表示と整合する定理であり、
\emph{graded algebra}の同型類を正確に記述したものである。
四頂点の任意周期付きZ代数の全同型を
$\operatorname{Aff}_3$で尽くしたとは主張しない。
四頂点potentialの商とgraded表示商を一般開上で同型にするには、
余分なZhang同一視を排除する別の比較が必要である。
\end{remark}

\begin{proposition}[八次臨界値多項式のアフィン共変性]
\label{prop:critical-discriminant-covariance}
$f\in\mathcal F_{19}^{\rm sm}$のJacobian環
$J_f=\kk[u_1,u_2,u_3]/(\partial_1f,\partial_2f,\partial_3f)$と
八次多項式$\Delta_f(\lambda)=\det(\lambda+m_f\mid J_f)$
を考える。
$f'=\rho_{A,b}(f)$とすると、アフィン変換$u\mapsto Au+b$は
$\operatorname{Spec}J_{f'}$と$\operatorname{Spec}J_f$を同型にし、
\begin{equation}\label{eq:critical-discriminant-covariance}
 \boxed{\quad
 \Delta_{f'}(\lambda)
 =(\det A)^{-8}\,
   \Delta_f\bigl((\det A)\lambda-f(b)\bigr)
 \quad}
\end{equation}
が成立する。これは非被約な臨界点schemeも含めて成り立つ。
特に、一般の8個のnodalな特異メンバーの対応は、
次数付き中心拡大$B_f$の同型に対して自然である。
\end{proposition}
\begin{proof}
式\eqref{eq:affine-action}を微分すると
$\nabla f'(u)=(\det A)^{-1}A^T\nabla f(Au+b)$となる。
従って二つのJacobianイデアルはアフィン座標変換で対応する。
対応する有限長8の環上では、$f'$の乗法は
$(m_f-f(b)\,\mathrm{id})/\det A$に移る。
その特性多項式を取れば式\eqref{eq:critical-discriminant-covariance}を得る。
この議論は固有値の相異性を使わず有限長環上で行うので
非被約な臨界点も含む。
一般の8つの通常二重点については
定理\ref{thm:pencil-critical}を適用する。
\end{proof}

\subsection{任意の周期を持つ関係式stackとの一般点での比較}
\begin{proposition}\label{prop:period-formal-rigidity}
$\check B$のFrobenius標準4周期点を取り、二次関係$R\subset B_1^{\otimes2}$の
二因子射影stabilizerが有限であるとする。
この点の全Artin標準AS四窓変形は、射影非捩れpotentialで一意に表示される。
従ってこの一般点では、任意周期付き関係式stackに潜在的な追加変形方向はない。
商に使うGLまたはPGLおよびスカラーrigidificationは両側で揃える。
\end{proposition}
\begin{proof}
Artin底$T$上の四窓のKoszul双対$E_T$は、長さ$0,\dots,4$で
$(1,4,6,4,1)$型の相対Frobenius代数である。
このことは標準AS条件の線型解消、または固定内部長の有限自由bar複体の底変更から従う。
最高次数の四つのtraceを選び、特殊ファイバーでは通常の意味で対称なtraceに規格化する。
そのNakayama自己同型$\nu_T$は各頂点を固定し、特殊ファイバーで恒等となる。
逆双対の次数一写像は隣接する二次関係を保存する。
小拡大の核$I$において$\nu=1+I D$と書くと、その射影像は特殊ファイバーの二因子
stabilizerのLie代数に入る。これは零なので、帰納的に各矢空間上の$\nu_T$はスカラー$c_i$である。
Nakayamaは最高次数traceを固定するから$\prod_{i=0}^3c_i=1$。
四つのtraceを単元$t_i$倍すると$c_i$は$(t_i/t_{i+1})c_i$へ変わり、積が1なので全て1にできる。
さらに頂点$i$のtraceを$(-1)^i$倍すればgraded対称traceとなる。
これを四次テンソルに双対化し、完全対合を用いれば関係空間はその二階微分となる。
二つのこのtraceは隣接完全対合により共通スカラーだけ異なるので、射影potentialは一意である。
この構成は関係式の同型と両立する。
\end{proof}
二因子stabilizerの有限性は第\ref{sec:surface-witness}節に記した$60\times32$行列のrank30によって、同じ$F_*$で直接確認される。
従って命題\ref{prop:period-counterexample}の大域的な不一致と、
$R(1,3)$一般点での形式的な一致は両立する。

\section{反標準環と認識部分のambient射}
次数一生成の次数付き環$A$に対し、その\emph{二次包絡}を
\[
 Q(A):=T_\kk(A_1)\Big/
 \big\langle\ker(A_1\otimes_\kk A_1\longrightarrow A_2)\big\rangle
\]
と定める。これは二次関係のデータを抽出する構成であり、
AS正則性は定義だけからは従わない。

\begin{proposition}[反標準環とambient候補]\label{prop:intrinsic-ring}
定義\ref{def:weak-cubic}に対し
\[
 \Gamma_{\rm ac}=\bigoplus_{m\ge0}\Hom(\OO,H^m\OO)
\]
は積も含めて内在的である。通常のArtin--Zhang再構成の有限性仮定の下で
$(\mathcal A,\OO,H)$は$(\qgr\Gamma_{\rm ac},\pi\Gamma_{\rm ac},(1))$から回復される。
$\Gamma_{\rm ac}$が次数一生成なら$Q(\Gamma_{\rm ac})$は一意な二次ambient候補である。
\end{proposition}
\begin{proof}
積は$f\in\Hom(\OO,H^m\OO)$、$g\in\Hom(\OO,H^n\OO)$に対して
$H^m(g)\circ f$である。Serre同値は特徴付けにより同値で保存され、
その同型の選択の変更はorbit algebraの同型を与える。
AZ amplenessと有限性の下でsection algebra再構成\cite{AZ}を適用する。
二次包絡は長さ二の合成核だけから定義される。
\end{proof}

内在的公理（rank-one torsion-free性とexceptionalityを含む）だけから、次数一生成や二次包絡のAS正則性が従うかは未証明であり、本稿では保留する。
本稿ではその一般的な反例を与えてはいないが、これらの含意は未証明である。
そこで内在的stackの\emph{超曲面認識部分関手}$\mathfrak C^{\rm hyp}$を、
$A=\Gamma_{\rm ac}$が次数一生成、$B=Q(A)$が標準AS4、
$I_3=\ker(B_3\to A_3)$がrank1でその生成元が左右正則、という条件で切る。
Hilbert級数と長さ4の核の次元4から左右乗法の像が一致し、
正規性と全高次核の生成が従って$A=B/(G)$となる。
族の局所自由性と分裂核から任意の底変更とも両立する。
従ってこの部分には実際の射
\begin{equation}\label{eq:ambient-map}
 \Phi:\mathfrak C^{\rm hyp}\longrightarrow\mathfrak P^{\rm AS,pt}_3,
 \qquad(\mathcal A,\OO)\longmapsto Q(\Gamma_{\rm ac})
\end{equation}
がある。ここでは源の自然同型と対応する頂点スカラーを処理したpointed targetを用いる。
rawな$\mathrm{GL}_4^4$商を用いる比較では、源にも同じ周期表示・スカラーのデータを加えた
presentation stackを用い、その正規incidenceの射を取る。内在的な源とrawなtargetを
そのまま比較して同じstack次元・同じファイバー次元を主張しない。
次数付きorbit algebraからの選択周期は次数shiftの4乗である。
非捩れpotentialのtargetを使う場合は適合周期へ規格化する。
この二つの選択は一般には同じ対象でなく、$R(1,3)$のCY表示上では適合する。

\section{AOUからの有限なambient構成と一般点の支配性}\label{sec:aou-ambient}
AOUの関係式族は、三つのblockを持つacyclic quiverの代数
$\Lambda$で、9頂点、18矢、9本の長さ2の関係を持つ。
長さ3の経路がなく、関係間のoverlapがないので、通常の相対quiver解消は長さ2で完全である。
$\Lambda$のrankは$9+18+18=45$である。
従って相対inverse Serre bimodule
$\Theta_\Lambda=\operatorname{RHom}_{\Lambda^e}(\Lambda,\Lambda^e)$を有限複体で計算できる。
$P_0$を頂点$(0,0)$の右射影対象とし、
\begin{equation}\label{eq:finite-serre}
 K_m=\operatorname{RHom}_\Lambda
       (P_0,P_0\mathbin{\otimes^L_\Lambda}\Theta_\Lambda^{\otimes^L m}[2m])
       \quad(0\le m\le4)
\end{equation}
と置く。これは係数環上の有限perfect複体である。

\begin{proposition}[低次数からの代数的候補射]\label{prop:aou-candidate}
滑らかな可換標識付き三次曲面を含み、相対inverse Serre複体の底変更が成立するAOUのある開部分上で、
$K_m$は次数0に集中し、そのrankは順に$1,4,10,19,31$となる。
$V=H^0K_1$、$R=\ker(V^{\otimes2}\to H^0K_2)$により、
有限Grassmannianの商への代数的射
\[
 \Psi_2:M_{\rm AOU}^{\rm low}\longrightarrow
 [\operatorname{Gr}(6,16)/\mathrm{PGL}_4]
\]
が得られる。これはまだ全域でAS性を結論する射ではない。
\end{proposition}
\begin{proof}
可換点では$P_0=\OO_X$、$\mathsf S^{-1}[2]=(-)\otimes\OO_X(1)$なので、
$K_m=R\Gamma(X,\OO_X(m))$である。$0\le m\le4$の高次コホモロジーは消える。
有限個のperfect複体に開性と底変更を適用し、rankと次数0集中を保つ開部分を取る。
$V^{\otimes2}\to H^0K_2$の全射性もこの点で成り立つので開に制限でき、核はrank6。
積はSerre orbitの合成から定まり、基底変更に両立する。従って所望の商への射がある。
\end{proof}


\begin{proposition}[一般幾何学的点でのinverse Serre orbit環]\label{prop:aou-orbit-generic}
AOU\cite[\S9, Theorem 9.4]{AOU}の一般位置条件
（滑らかな楕円曲線、無限位数のtranslation、十分一般の六点）を満たす
楕円nonupleに付随する九頂点有限代数を$\Lambda$、
頂点$(0,0)$の例外対象を$P_0$とし、
$H=\mathsf S^{-1}[2]$とする。積を保つ次数付き環同型
\[
 \bigoplus_{m\ge0}\Hom_{D^b(\Lambda)}(P_0,H^mP_0)
 \ \simeq\ F
\]
が存在する。ここで$F=P/(G)$はVan den Bergh\cite[Chapter 12]{VdB}の
反標準環であり、$P$は中心一次元$t$を持つ四次元AS正則Koszul代数、
$G$は正則正規三次元である。
この主張は幾何学的点での導来圏内の同型であって、
アーベルheartの一致や相対底変更を含まない。
\end{proposition}
\begin{proof}
AOU (9.44)--(9.45)の例外列の先頭は構造対象$\OO$であり、
(9.55)--(9.57)の導来同値は対応する$P_0$を
$\pi F\in\qgr F$へ送る。同論文Lemma 9.2により
$D^b\qgr F$のSerre関手は$(-1)[2]$なので、
$H$は次数shift$(1)$に対応する。
したがって左辺は
$\bigoplus_{m\ge0}\Hom_{\qgr F}(\pi F,\pi F(m))$
と積を保って同型になる。
$P$が四次元AS正則で$G$が正則三次元であることから、
$0\to P(-3)\xrightarrow{\cdot G}P\to F\to0$によって
$F$のgraded depthは少なくとも3である。
従って$H^0_{\mathfrak m}(F)=H^1_{\mathfrak m}(F)=0$であり、
Artin--Zhangのsection環の完全列\cite{AZ}を用いれば
$\Hom_{\qgr F}(\pi F,\pi F(m))\simeq F_m$が$m\ge0$で成立する。
\end{proof}

\paragraph{generic pointに用いる既知の入力.}
AOU\cite[Theorem 1.1, \S9, Theorem 9.4]{AOU}と
Okawa--Ueda\cite[Theorem 6.4]{OU2025}は、一般の標識付きデータを楕円曲線と六点blow-upで記述する。
命題\ref{prop:aou-orbit-generic}によりAOUの導来同値とSerre関手から、
一般幾何学的点でのinverse Serre orbit環は$F$に一致する。
同論文 (9.6)--(9.10) とVan den Bergh\cite[Chapter 12]{VdB}によれば
$F=P/(G)$で、$P$は中心一次拡大AS4、$G$は正則正規三次元。
幾何学的点での同型は命題\ref{prop:aou-orbit-generic}で証明したが、
相対底変更・形式変形への比較は残っている。
全標識付き成分やコンパクト化全体には拡張しない。

一般のSklyanin頂部では三次中心が一次元である。
正規化自己同型を$Gx=\sigma(x)G$と書くと、中心正則一次元$t$から$\sigma(t)=t$。
$P/(t)$上で$\overline G$は中心正則なので$\overline\sigma=1$となる。
よって$\sigma(x_i)=x_i+c_i t$である。
Pym/LBSの一般Jacobian座標で二次関係を比較すると
$\sum_i c_i\partial_i\partial_j f_3=0$であり、滑らかな$f_3$から$c_i=0$。
従ってこの一般部では$G$は中心的である。
標準の中心持上げと低次数中心の消滅により$Z(P)_3=\kk\Omega\oplus\kk t^3$となる。
一般の中心拡大の型がPymの$R(1,3)$であることは\cite[\S3.4]{Pym},\cite{LBS}の入力である。

\begin{theorem}[一般AOU点からのambientの像]\label{thm:aou-dominance}
命題\ref{prop:aou-orbit-generic}の一般位置に制限する。
さらに(i)この幾何学的点での復元が相対inverse Serre orbit環の底変更と両立し、
(ii)可換標識付き点からの形式変形の8方向をAOUの関係式表示がすべて実現し、
低次数写像の微分比較が適用できる、と仮定する。
この比較が成立する一般開部分では命題\ref{prop:aou-candidate}は
AS4 ambientへの有理射を与え、$R(1,3)$の7次元成分を支配する。
対応する中心三次対への有理像は8次元であり、その一般ファイバーは有限である。
従ってambientへの有理射の一般ファイバーは1次元である。
その有限次数を$|W(E_6)|$と同定したわけではない。
\end{theorem}
\begin{proof}
命題\ref{prop:aou-orbit-generic}により一般幾何学的点のorbit環は$P/(G)$であるから、二次包絡は$P$に一致する。
従って候補射の一般値は$R(1,3)$に属し、三次核もrank1となる。
これらは有限な式とrank条件なので、低次数の有理射として定義される。
支配性を以下で示した後、定理\ref{thm:r13-open}のatlasの逆像へ制限する。

像の次元を単なるパラメータ差から推測せず、truncationの微分を使う。
可換点においてAOUの8方向が全Abel変形を与えることは、定理の仮定(ii)の内容である。
反標準環の次数$0,1,2,3$の変形が、基底変更を除いて自明なら、その全次数変形も自明である。
実際、Artin底で次数一生成性は各次数のNakayamaにより持ち上がる。
自由テンソル代数からの核$J$は各次数で平坦であり、
$J/(VJ+JV)$の次数$>3$は特殊化で零なのでNakayamaにより零である。
従って核は次数2と3で生成され、低次数の自明化は環全体を自明化する。
AZ再構成によってAbel変形も自明となる。
よって低次数射の微分は単射であり、曲面同型類の像の次元は8である。
$A_1$の射影基底を加えると源の次元は$8+15=23$で、低次数像も23次元となる。

一般点ではこの低次数像は中心三次対であり、固定ambientに対する中心直線は$\PP^1$。
従ってambient像の次元は少なくとも$23-1=22$。
一方、中心拡大型ambientの座標付き像は22次元であるから等号。
一般の射影自己同型が有限である範囲で$\mathrm{PGL}_4$を割り、次元7と1次元一般ファイバーを得る。
四頂点potential側への移行は定理\ref{thm:r13-open}および
命題\ref{prop:period-formal-rigidity}の一般開部分を用いる。
\end{proof}

\begin{remark}[有理射と弱い対象の全域の射]
式\eqref{eq:finite-serre}は有限な行列・tensor積・核の計算による明示的レシピであり、
8個のトーラス座標から7個の正規形座標への展開済み有理式はまだ計算していない。
一般点の$\qgr F$はAOUの完全強例外列との導来同値によりsmooth properであり、
Serre公式とAZの定理により定義\ref{def:weak-cubic}を満たす。
従って式\eqref{eq:ambient-map}の像はこの成分の幾何学的一般点を含み、
その閉包の余次元は0である。
これは認識部分における\emph{一般点の支配性}であり、AOUの全開集合がそのまま
相対smooth properな族を与えることや、全境界を含む正則射、全内在的成分での認識定理は含まない。
\end{remark}


\section{三次曲面のPoisson構造の一意延長}
$S=\kk[x,y,z,t]$、$X=V_+(F)\subset\PP^3$を滑らかな三次曲面とし、
$\Omega=dx\wedge dy\wedge dz\wedge dt$を固定する。
随伴公式より$-K_X=\OO_X(1)$なので
\[
 H^0(X,\wedge^2T_X)=H^0(X,\OO_X(1))\simeq S_1.
\]
次元2では全双ベクトルがPoissonである。
$F$の代表と$\Omega$を固定して、この同定の規格化を選ぶ。

\begin{theorem}\label{thm:extension}
任意の$h\in S_1$について
\[
 \iota_{\Pi_{F,h}}\Omega=dF\wedge dh
\]
で定まる斉次二次Poisson双ベクトルは$\PP^3$へ降り、$X$をPoisson部分多様体とし、
$X$上の$h$に対応するPoisson構造を誘導する。
$X$上の任意のPoisson構造の$\PP^3$への延長は一意である。
\end{theorem}
\begin{proof}
$h\ne0$なら線型座標を変えて$h=t$とできる。このとき$t$はCasimirであり、
\[
 \{x,y\}=F_z,\quad\{y,z\}=F_x,\quad\{z,x\}=F_y.
\]
$t$を係数とする三変数Jacobian括弧なのでJacobi恒等式を満たし、$F$もCasimirである。
斉次次数から次数零の有理式上の括弧が定まり、射影化できる。
$t=1$上での随伴の局所式を使えば、誘導双ベクトルは反標準切断$t$に対応する。
$h=0$の場合は零である。

二つの延長の差は
$H^0(\PP^3,\wedge^2T_{\PP^3}(-3))$に属する。一方
\[
 \wedge^2T_{\PP^3}(-3)\simeq\Omega^1_{\PP^3}(1),\qquad
 H^0(\PP^3,\Omega^1_{\PP^3}(1))=0
\]
であり、最後の消滅はEuler完全列から従う。
よって延長は一意である。
\end{proof}

\begin{theorem}\label{thm:poisson-fiber}
非零Poisson構造を持つ滑らかな三次曲面のincidenceからambient Poisson構造への射は、
$R(1,3)$の錐の稠密部分を覆う。
一般の固定ambient上のファイバーは
\[
 X_c=V_+(F+c h^3),\qquad c\in\mathbb A^1
\]
のうち滑らかな曲面を与える開部分である。Poissonスカラーを残す一般有理商の次元は
\[
 8\longrightarrow7,
\]
一般ファイバーの次元は1である。
\end{theorem}
\begin{proof}
$h=t$として最高active次数部分$f=F(x,y,z,0)$が滑らかな範囲を取る。
この括弧の一次Poisson中心は$\kk t$である。
実際、$a x+b y+c z$が中心なら定数ベクトル$(a,b,c)$と$\nabla f$は常に平行となり、
$f$が一つの線型形式の三乗となって滑らかさに矛盾する。
従って同じ非零ambientを与える$h'$は$h$の定数倍である。
$h$を揃えた後、$d(F'-F)\wedge dh=0$から
$F'-F\in\kk[h]_3=\kk h^3$となる。
逆にこの変換は括弧を変えない。ペンシルの残る端点$h^3$は三重平面である。
滑らかさは開条件であり、最初の$F$が滑らかならこの開部分は非空。

ambientのパラメータは中心一次直線$[h]\in\PP(S_1)$の3次元と、
$F\bmod\kk h^3$の19次元で合計22次元である。
基底変更の下で係数は捩れるが、局所chartでの次元計算には影響しない。
像の既約閉包は$R(1,3)$の錐である\cite[\S3.4]{Pym}。
曲面とPoisson構造のincidenceは$\PP(S_3)^{\rm sm}$上の階数4の束で23次元。
一般の$\PGL_4$軌道は15次元であることを次節で厳密に確認する。
従って一般有理商は$23-15=8$、$22-15=7$である。
固定同型類のファイバーには有限自己同型群による商があり得る。
\end{proof}

\begin{remark}
固定ambientは、有理第一積分$F/h^3$の値を忘れている。
三次曲面はその一般葉の閉包であり、ペンシル全体で楕円曲線
$V(F,h)$を共有する。これは幾何学的な1次元忘却の説明である。
$\Pi$自体をスカラー射影化するとambientの次元は21になる。
本稿の22と7はPoissonスカラーを残す規約である。
\end{remark}


\begin{theorem}[中心三次ペンシルの臨界値多項式]\label{thm:pencil-critical}
$F\in\kk[x,y,z,t]_3$について$f_3:=F(x,y,z,0)$が滑らかな
平面三次曲線を定めるとし、$f:=F(x,y,z,1)$と置く。
\[
 J_f=\kk[x,y,z]/(f_x,f_y,f_z),\qquad
 \Delta_f(\lambda)=\det(\lambda\,\mathrm{id}_{J_f}+m_f)
\]
と定める。ここで$m_f$は$f$による乗法である。
このとき$\dim_\kk J_f=8$、$\Delta_f$はモニック8次であり、
$X_\lambda:=V_+(F+\lambda t^3)\subset\PP^3$が特異であることと
$\Delta_f(\lambda)=0$は同値である。
さらに係数空間の非空Zariski開集合上で$\Delta_f$は8個の相異なる根を持ち、
各特異$X_\lambda$にはただ一つの通常二重点がある。
この開集合では、滑らかな三次曲面のパラメータは
$\mathbb A^1\setminus V(\Delta_f)$である。
\end{theorem}
\begin{proof}
三つの二次式$f_x,f_y,f_z$を斉次化して得る$\PP^3$上の共通零点は
無限遠平面$t=0$に存在しない。実際そのような点では
$(f_3)_x=(f_3)_y=(f_3)_z=0$となり、
Eulerの恒等式と$f_3=0$の滑らかさに矛盾する。
従って三つの二次式は零次元完全交差を定め、
B\'ezoutの定理によって長さは$2^3=8$である。
全ての零点は$t\ne0$に属するから$\dim_\kk J_f=8$となる。

$X_\lambda$は$t=0$に特異点を持たない。
$t=1$では方程式が$f+\lambda=0$となり、
特異点の条件は$f_x=f_y=f_z=f+\lambda=0$である
（$t$偏微分の消滅はEulerの恒等式による）。
従って$\lambda$は$m_f$の固有値の負であることが必要十分である。

非空性には$F_0=x^3+y^3+z^3-t^2(3x+12y+48z)$を考える。
$\nabla f_0=0$の8点は$(\pm1,\pm2,\pm4)$であり、
Hessianは$\operatorname{diag}(6x,6y,6z)$だから全て非退化である。
これらの点での臨界値は
$\{-146,-142,-114,-110,110,114,142,146\}$で相異なる。
よって$\Delta_{f_0}$の判別式は非零である。
この条件およびJacobian環のエタール性は開条件だから、
同じ性質は非空Zariski開部分で成立する。
各特異曲面には非退化臨界点が一つだけあり、通常二重点となる。
\end{proof}
\begin{remark}
$\lambda=\infty$の端点は三重平面$t^3=0$であり、
上の8個のnodalな曲面には数えない。
また、この定理は可換Poissonペンシルに関するものであり、
量子化した$B/(\Omega+\lambda t^3)$の特異値も同じであるとは主張しない。
有理数上の検算は\texttt{computations/verify\_pencil\_discriminant.py}にある。
\end{remark}


\begin{proposition}[八点の臨界集合と非可換point module]\label{prop:critical-point-modules}
命題\ref{prop:potential-presentation}のJacobian型中心拡大$B_f$を考え、
$F(x,y,z,t)=t^3 f(x/t,y/t,z/t)$と書く（$t^3$の定数項は任意）。
すると、$t$が非零に作用する$B_f$のpoint moduleのパラメータschemeは
\[
 \operatorname{Spec} J_f,\qquad
 J_f=\kk[x,y,z]/(f_x,f_y,f_z)
\]
と自然に同型である。従って定理\ref{thm:pencil-critical}の一般開部分では
この型のpoint moduleはちょうど8個存在し、
それぞれがペンシルの一つのnodalな特異曲面を指定する。
この対応は係数$f$に依存するもので、量子中心三次元の固有値を
可換臨界値と同定する主張ではない。
\end{proposition}
\begin{proof}
$t$は中心正則な次数1元であるから、$t$が各次数で非零に作用する
point moduleは、基底を$t\,m_n=m_{n+1}$となるよう選べる。
$t x_i=x_i t$により$x_i\,m_n=a_i m_{n+1}$の係数$a_i$は
$n$に依存しない。命題\ref{prop:potential-presentation}の二次関係
$[y,z]=\operatorname{sym}F_x$とその巡回式をこのmoduleに作用させれば、
左辺は零となり、右辺はそれぞれ$f_x(a)$、$f_y(a)$、$f_z(a)$となる。
逆に$\nabla f(a)=0$なら上の作用は全ての関係を満たし、
point moduleを与える。scheme構造についても、
$t$を逆元化した次数零部分$B_f[t^{-1}]_0$の一次元表現環は
交換子関係の可換化で得られる$J_f$である。
最後の結論は定理\ref{thm:pencil-critical}から従う。
\end{proof}

\section{Poisson成分の一般正則性と有限行列検証}\label{sec:surface-witness}
次の有理係数の点を取る。
\begin{align*}
 F_*={}&3(x^3+y^3+z^3)+5xyz+t(x^2+2y^2+3z^2)\\
 &+t^2(x+2y+4z)+7t^3,\qquad h=t.
\end{align*}
勾配イデアルが零次元であることから$V_+(F_*)$は滑らかである。
この確認は有理数上のGr\"obner計算で再現する。
二次双ベクトルの係数空間は60次元。
全Jacobi方程式80係数と発散零16係数の微分を$J$とし、
$F$の20係数と$h$の4係数の変分を$M$、$\mathfrak{gl}_4$の軌道微分を$O$とする。
付属コードは、整数上の$JM=0$と次の非零有理小行列式を検証する。
\[
\begin{array}{c|r|r|l}
 &\text{行列の大きさ}&\rank&\text{選択した最大非零minor}\\\hline
M&60\times24&22&41472\\
O&60\times16&15&407413613894400\\
J&96\times60&38&483917294055304538158847754240000\\
H&60\times32&30&-12020528624208926290875/2
\end{array}
\]
行・列の添字は\texttt{surface\_incidence\_certificate.json}に固定する。
$H$については下記で説明する。

\begin{proposition}\label{prop:reduced}
$R(1,3)$の22次元錐は全unimodular二次Poissonスキームの既約成分であり、
一般点で正則である。
\end{proposition}
\begin{proof}
$\rank M=22$によりJacobian族の像の局所次元は22以上であり、
定理\ref{thm:poisson-fiber}のパラメータ数から22以下である。
$\rank J=38$なので全スキームの接空間次元も$60-38=22$。
従ってこの点の局所環は正則であり、22次元の既約像の閉包はその成分となる。
\end{proof}

$\Pi=\Pi_{F,t}$に対する二次graph関係を
\[
 R_\Pi=\langle x_i\otimes x_j-x_j\otimes x_i-\operatorname{sym}\Pi_{ij}\rangle
\]
とする。$H$は$(g\otimes h)R_\Pi=R_\Pi$を恒等で線型化し、
$(V\otimes V)/R_\Pi\simeq\Sym^2V$へ写した行列である。
二つの独立な矢空間スカラーが核に入り、rank30によりこれが全核である。
従って$(\PGL_4)^2$内の二因子stabilizerは、この点および一般点で有限である。
この有限検査自体は、当該代数の全AS正則性を証明するものではない。
Pymの中心拡大族の既知の良い範囲とこの非零minor開集合を交差させて使う。

\section{量子側の復元と一次元ファイバー}\label{sec:surface-fiber}
まず$B=T(V)/(R_2)$は二次、$G\in B_3$は正則正規元、$A=B/(G)$とする。
三次商は次数一・二を変えないので
\begin{equation}\label{eq:envelope}
 B=\frac{T(A_1)}{\langle\ker(A_1\otimes A_1\to A_2)\rangle},\qquad
 \ker(B_3\to A_3)=\kk G.
\end{equation}
これは次数付き同型について関手的であり、相対的に成分が局所自由なら二次核の形成は
任意の底変更と可換する。存在している二次ambientの一意性と、任意の圏からのambientの
存在は別の主張である。

\begin{proposition}[中心三次ペンシル]\label{prop:center}
$B$は中心正則一次元$t$を持ち、active次数による濾過が
$\operatorname{gr}B\simeq S_0\otimes\kk[t]$を満たすとする。
ここで$S_0$は三変数AS正則整域で
\[
 Z(S_0)_1=Z(S_0)_2=0,\qquad Z(S_0)_3=\kk\omega.
\]
さらに$\omega$を最高active記号に持つ中心三次元$\Omega_B$が存在すると仮定する。
このとき
\[
 Z(B)_1=\kk t,\qquad Z(B)_3=\kk\Omega_B\oplus\kk t^3.
\]
固定ambientへの同定を保持した中心三次商の空間は$\PP^1$であり、
$[t^3]$を除くchartは$B/(\Omega_B+\lambda t^3)$である。
有限な射影自己同型群で割っても一般ファイバー次元は1である。
\end{proposition}
\begin{proof}
中心元の最高active記号も中心的である。三次元のactive次数が3なら
$\Omega_B$の定数倍を引き、残るactive次数2と1の記号は仮定によって零となる。
最後は$t^3$の定数倍だけが残る。一次の場合も同様。
随伴次数付き代数が整域なので$B$は整域であり、非零中心元は正則である。
式\eqref{eq:envelope}により商の次数付き同型はambientの自己同型へ持ち上がる。
従って固定同型類のファイバーは自己同型群による商となる。
\end{proof}

この仮定を満たす中心拡大と中心元は\cite{LBS,EG,Pym}の構成に属する。
ここでは無限個の次数すべてで中心が一定な一つの開集合は要求しない。
次数3までの条件だけを使っている。
一般の良いgraph族で、中心一次直線の3係数と三次potentialの19係数からambient22次元、
中心三次ペンシルを加えて23次元となる。有限射影stabilizerの範囲で一般有理商は
7次元と8次元である。これは指定した次数付き中心拡大族についての主張であり、
AOUの全コンパクト化との双有理同定を前提としない。
定理\ref{thm:aou-dominance}が与える一般AOUからの支配性は、この結論と低次数復元を併用する。

\begin{proposition}[固定ambientの全正規形式ファイバー]\label{prop:normal}
命題\ref{prop:center}の$B$について、二因子stabilizer
$H_R\subset\PGL(B_1)^2$が有限であるとする。
このとき中心三次元$G$における、
固定した$\check B$上のindexed正規三次イデアルの形式変形は、
中心三次直線の形式変形に一致する。従って厳密形式ファイバーは
$\operatorname{Spf}\kk[[s]]$である。
\end{proposition}
\begin{proof}
局所Artin底$R$上で、各三次核直線の生成元$G_i$を$G$へ還元するように選ぶ。
正規性と正則性から$G_i x=\varphi_i(x)G_{i+1}$となる可逆線型写像を得る。
二次関係へ適用し、右から$G_{i+2}$を消去すると
$(\varphi_i\otimes\varphi_{i+1})(R_2)=R_2$である。
特殊ファイバーでは射影的に恒等である。標数零の有限代数群$H_R$はエタールなので、
Artin底で恒等を持ち上げる射影対は恒等だけである。
従って各$\varphi_i$はスカラーであり、$G_i$のスカラーを添字ごとに再帰的に変更して
$G_i x=xG_{i+1}$にできる。$x=t$と置けば$t$は中心正則なので$G_i=G_{i+1}$。
よって一つの中心元を全添字で反復した族となる。
逆に中心三次直線はそのような族を与える。
ここで$B$は次数一生成なので、$Z(B)_3$は有限次元線型写像
\[
 B_3\longrightarrow\Hom_\kk(B_1,B_4),\qquad
 u\longmapsto(v\mapsto uv-vu)
\]
の核である。体$\kk$上の任意の局所Artin代数$T$は平坦なので、
この核は$T$へのテンソル積と可換する。従って
$Z(B\otimes_\kk T)_3=Z(B)_3\otimes_\kk T$となる。
命題\ref{prop:center}から、この中心空間はArtin底でもrank2である。
従って$\PP^1$の一点での完備化である。
\end{proof}

\section{反標準偏極による内在性と局所変形}
三次曲面の場合、三次三様体と異なり$-K_X=\OO_X(1)$である。
非可換側でも、smooth proper性とsaturationを持つ中心CY型三次商について
Serre関手が$\mathsf S_A=(-1)[2]$なら
\[
 H=\mathsf S_A^{-1}[2]=(1),\qquad
 A_n=\Hom(\OO_A,H^n\OO_A)
\]
と復元できる。従って$(D^b\qgr A,\OO_A)$から$A$と$B=Q(A)$を復元でき、
偏極の平方根を選ぶ問題はない。上記Serre公式はAOU\cite[Lemma 9.2]{AOU}でも
所定のblow-up表示について確認されている。
任意の正規商では中山自己同型の捩れがあり得るため、上の公式を無条件には使わない。
構造対象自体を忘れた導来圏からの復元もここでは主張しない。

\begin{proposition}\label{prop:hh}
滑らかな可換三次曲面$X$について
\[
 \HH^1(X)=0,\qquad \dim\HH^2(X)=4+4=8,\qquad \HH^3(X)=0.
\]
従ってHochschild複体で制御される通常のmarked形式変形は無障害で、
そのminiversalな底は8次元形式円板である。
\end{proposition}
\begin{proof}
HKRにより$\HH^2=H^1(T_X)\oplus H^0(-K_X)\oplus H^2(\OO_X)$。
$H^2(\OO_X)=0$、$h^0(-K_X)=4$である。
接束の正常完全列とEuler完全列から$H^0(T_X)=0$および
$h^1(T_X)=h^0(\OO_X(3))-15=19-15=4$を得る。
前者の消滅は、$F$の偏導関数が二次正則列であり、一次係数のsyzygyがないことにもよる。
$H^2(T_X)=0$は正常完全列の高次コホモロジー、
$H^1(-K_X)=0$は$X$の直線束の中間コホモロジー消滅から従う。
$\wedge^3T_X=0$なので$\HH^3=0$。
同様に$\HH^1=H^0(T_X)\oplus H^1(\OO_X)=0$。
通常のHochschild変形理論では障害は$\HH^3$に入るため無障害性が従う。
この計算は反標準環の全域AS性や量子化の代数化を証明するものではない。
\end{proof}

\section{二種類の箙と、未完の大域比較}
AOUのmarked曲面模型では9頂点18矢の箙に9本の二次関係があり、
3-CY completionは9頂点27矢の三次potentialを持つ。
非可換$\PP^3$の標準4周期模型では4頂点16矢、四次superpotentialの二階微分を使う。
前者は曲面の例外列・helixに付随し、後者は4次元AS正則ambientに付随する。
CY次元も3と4で異なり、同じ箙の単純な部分potentialではない。
曲面の反標準環自体はAS Gorenstein次元3・index1であり、
9頂点のAS正則代数ともambientのAS正則代数とも区別する。

比較の具体的手順は、AOUの有限代数$\Lambda$と指定頂点対象$P_0$から
\[
 A_m=\Hom_{D^b(\Lambda)}
       (P_0,(\mathsf S_\Lambda^{-1}[2])^m P_0)
\]
を作り、低次数での積を計算して$Q(A)$を得ることである。
これは構造対象とSerre関手が同定された良い範囲でのレシピである。
全体で射にするには、相対Homの底変更、豊富性、二次包絡のAS性、正規三次核、
heartとmarkingの同値関係を証明する必要がある。

\paragraph{現在の結論と限界.}
定理\ref{thm:r13-open}は四頂点potentialの全方向で成分の一般正則性と7次元atlasを与える。
命題\ref{prop:period-formal-rigidity}はこの一般点で任意の選択周期の追加方向を排除するが、
命題\ref{prop:period-counterexample}が示す大域的な相違は残る。
定理\ref{thm:formal-weyl}は可換点の形式Weyl商であり、
大域商のためには命題\ref{prop:global-weyl-criterion}のmarking認識を証明する必要がある。
定理\ref{thm:aou-dominance}では一般AOUからの有理ambient射と余次元0を得たが、
全weakly marked対象が超曲面認識部分に入ること、コンパクト化全体で射が正則であること、
全境界でのsmooth proper性は未証明である。
全potential境界の明示分類、7係数atlasの有限同一視の次数、Weyl商の大域的な一般次数も未決定である。

\paragraph{検証と再現.}
\texttt{verify\_r13\_potential.py}は19独立係数を保つ微分恒等式、
四flatteningのrank、全接空間、群作用、7係数横断性を有理数上で検証する。
\texttt{verify\_surface\_incidence.py}はPoisson側の独立な接空間と二因子stabilizerを再現する。
有限行列計算と、本文の全次数のAS性・圏の降下の証明は区別する。
旧版に対する計算と組版の記録は\texttt{research/surface\_validation.json}にある。
本改訂版に対する再実行・外部査読・形式検証・新規性の網羅的確認は未完了である。

\begin{thebibliography}{99}
\bibitem{AOU} T. Abdelgadir, S. Okawa, K. Ueda,
\emph{Compact moduli of marked noncommutative cubic surfaces},
\href{https://arxiv.org/abs/2404.00175}{arXiv:2404.00175}, v1.
参照: Introduction, \S5.4, Theorems 1.1, 9.4, Lemma 9.2.
\bibitem{VdB} M. Van den Bergh,
\emph{Blowing up of non-commutative smooth surfaces},
Memoirs Amer. Math. Soc. 154 (2001), no. 734;
\href{https://arxiv.org/abs/math/9809116}{arXiv:math/9809116}, Chapter 12.
\bibitem{Pym} B. Pym, \emph{Quantum deformations of projective three-space},
Adv. Math. 281 (2015), 1216--1241; \href{https://doi.org/10.1016/j.aim.2015.06.005}{doi:10.1016/j.aim.2015.06.005};
\href{https://arxiv.org/abs/1403.6444}{arXiv:1403.6444}, \S2, \S3.4.
\bibitem{LBS} L. Le Bruyn, S. P. Smith, M. Van den Bergh,
\emph{Central extensions of three-dimensional Artin--Schelter regular algebras},
Math. Z. 222 (1996), 171--212; \href{https://doi.org/10.1007/PL00004532}{doi:10.1007/PL00004532}.
\bibitem{EG} P. Etingof, V. Ginzburg,
\emph{Noncommutative del Pezzo surfaces and Calabi--Yau algebras},
J. Eur. Math. Soc. 12 (2010), 1371--1416;
\href{https://arxiv.org/abs/0709.3593}{arXiv:0709.3593v4}, Theorem 3.3.2, \S9.

\bibitem{BSW} R. Bocklandt, T. Schedler, M. Wemyss,
\emph{Superpotentials and higher order derivations},
\href{https://arxiv.org/abs/0802.0162}{arXiv:0802.0162v3}, Definitions in \S1 and Theorems 6.2, 6.8.
\bibitem{MS} I. Mori, S. P. Smith,
\emph{The classification of 3-Calabi--Yau algebras with 3 generators and 3 quadratic relations},
\href{https://arxiv.org/abs/1502.07403}{arXiv:1502.07403v2}, Theorem 4.2(1).
\bibitem{AZ} M. Artin, J. J. Zhang,
\emph{Noncommutative projective schemes}, Adv. Math. 109 (1994), 228--287,
\href{https://doi.org/10.1006/aima.1994.1087}{doi:10.1006/aima.1994.1087}.
\bibitem{Lowen} W. Lowen,
\emph{Obstruction theory for objects in abelian and derived categories},
Comm. Algebra 33 (2005), 3195--3223,
\href{https://arxiv.org/abs/math/0407019}{arXiv:math/0407019}, Theorem A.
\bibitem{DDL} O. De Deken, W. Lowen,
\emph{Abelian and derived deformations in the presence of $\mathbb Z$-generating geometric helices},
J. Noncommut. Geom. 5 (2011), 477--505,
\href{https://doi.org/10.4171/JNCG/83}{doi:10.4171/JNCG/83},
Theorems 4.5, 5.13, 5.15 and Definition 5.12.
\bibitem{OU2025} S. Okawa, K. Ueda,
\emph{Marked Artin--Schelter surfaces of del Pezzo types},
\href{https://arxiv.org/abs/2508.04986}{arXiv:2508.04986v1} (2025),
Definition 1.1 and Theorem 6.4.
\bibitem{OWR} K. Ueda (joint work with S. Okawa),
\emph{Affine Weyl group actions on moduli spaces of marked noncommutative del Pezzo surfaces},
Oberwolfach Reports 44/2024, pp. 2557--2558,
\href{https://doi.org/10.4171/OWR/2024/44}{doi:10.4171/OWR/2024/44}, Conjectures 1--3.

\end{thebibliography}
\end{document}
